% Imporditud: tools/import_eksperimendid.py
% Allikas: Eesti füüsikaolümpiaadi eksperimendi ülesannete
%   kogu (Rael Kalda, 2023), ülesanne Ü117, lahendus L117.
\setAuthor{Mihkel Kree}
\setRound{piirkonnavoor}
\setYear{2017}
\setNumber{G E1}
\setDifficulty{3}
\setTopic{dünaamika}
\setVahendid{statiiv, nöör, koormis, joonlaud, stopper, millimeeterpaber}
\setPunktid{10}
\setAllikas{Eksperimendi kogu Ü117, lahendus lk 92}
\setLahendusseis{kasitsi_tehtud}

\prob{Pendlikonstant}
On teada, et pendli võnkeperiood $T$ avaldub kujul $T = \alpha\sqrt{l}$, kus
$l$ on pendli pikkus.

Määrake võrdeteguri $\alpha$ väärtus, kasutades oma mõõtetulemuste graafilist
esitust.

\hint

\solu
% Käsitsi ümber kirjutatud kogumiku lahendusest L117 (lk 96). Tabel ja
% hindamisskeem on PDF-i tekstikihist loetamatud, seega on nad siin uuesti
% laotud. Arvväärtused on kogumiku omad ja üle kontrollitud:
% T = 20T/20 ning alpha = sqrt(tous), tous 0,0406 s^2/cm -> alpha = 0,201.
Mõte on teha graafik \emph{sirgeks}. Sõltuvus $T = \alpha\sqrt{l}$ ei ole
sirge, aga kui kanda teljestikku $T^2$ sõltuvana $l$-st, siis kehtib
\[
T^2 = \alpha^2 l
\]
ja tõusust saab $\alpha$ kätte. Sama hästi võib panna teljele $\sqrt{l}$ ja
$T$. Sirge sobitamine katsepunktide vahele on palju kindlam kui ruutjuure
kuju sobitamine.

Aja määramatuse vähendamiseks tuleb mõõta korraga võimalikult palju
perioode: alg- ja lõpphetke fikseerimise viga jaguneb siis kõigi võngete
peale ära. Pendlile antakse väike hoog, sest suure amplituudi korral ei ole
periood enam amplituudist sõltumatu. Pikk nöör on parem, sest siis ei mõjuta
koormise enda inertsimoment ja massikeskme asukoha teadmatus tulemust nii
palju. Pika nööri mõõtmiseks tehakse nöörile sõlmed või märgid, mille
vahekaugused mõõdetakse joonlauaga ja liidetakse kokku.

Kogumiku katses on viis pikkust ja iga kord mõõdetud 20 perioodi:

\begin{center}
\begin{tabular}{lccccc}
\hline
$l$ / cm & 50 & 60 & 70 & 80 & 90 \\
$20T$ / s & 28,04 & 31,38 & 33,62 & 36,01 & 38,23 \\
$T$ / s & 1,40 & 1,57 & 1,68 & 1,80 & 1,91 \\
$T^2$ / s$^2$ & 1,97 & 2,46 & 2,83 & 3,24 & 3,65 \\
\hline
\end{tabular}
\end{center}

Sobitatud sirge tõus on $\SI{0,0406}{s^2/cm}$ ja kuna tõus on $\alpha^2$,
\[
\alpha = \sqrt{\num{0,0406}} = \SI{0,201}{s/cm^{0,5}}
= \SI{2,01}{s/m^{0,5}}
\]
Teooria annab $\alpha = 2\pi/\sqrt{g} \approx \SI{2,02}{s/m^{0,5}}$, seega
sama katse annab ka raskuskiirenduse.

\textbf{Hindamisskeem} (kokku 10 p): katse kirjeldus \pp{1} (oluline on ka
dokumenteerida, mida täpselt tehakse); mõõtmised dokumenteeritud \pp{1};
perioodi mõõtmine vähemalt viie erineva nööri pikkusega \pp{2}, vähemalt
kolme pikkusega \pp{1}; korraga vähemalt kümne perioodi mõõtmine \pp{2},
vähemalt viie perioodi mõõtmine \pp{1}; korrektne graafik koos sirgega
\pp{2}; graafikult leitud vastus ei erine teoreetilisest väärtusest rohkem
kui \SI{10}{\percent} \pp{2}, rohkem kui \SI{20}{\percent} \pp{1}. Kui
õpilane leiab vastuse graafikut kasutamata, anda \pp{1} juhul, kui vastus ei
erine teooriast rohkem kui \SI{10}{\percent}.
\probend
